LiDAR range-image compression for autonomous driving is investigated as a staged, end-to-end learning problem with explicit ROI-aware quantization control. The implemented pipeline encodes 5-channel range images, predicts spatial importance, applies adaptive latent quantization, and reconstructs the input through a decoder. Four operational stages are defined in the current codebase: Stage0 (uniform baseline), Stage1 (adaptive quantization with ROI supervision), Stage2 (Stage1 plus teacher distillation), and Stage3 (Stage2 objective with multi-scale ROI head ablation).

This thesis draft formalizes the exact training objective used in code, including reconstruction loss, normalized rate proxy terms, weighted ROI BCE supervision, ROI/background separation margin, and composite distillation loss. In 150-epoch sweeps with balanced objectives and a higher-capacity ROI head, the best Stage1 run reached final loss 0.7224 (ResNet, lr=1e-4), while the best Stage2 run in the evaluated DarkNet distillation sweep reached 2.1185 at lambda\_distill=0.1. Short fair-budget pilots with native versus oracle ROI control showed a key diagnostic result: with oracle ROI, entropy bitrate proxy dropped substantially (e.g., Stage1: 0.8064 to 0.4177; Stage2: 0.8687 to 0.4492 in bpp-entropy) at nearly unchanged ROI-MSE.

These findings indicate that adaptive quantization can deliver strong bitrate gains when ROI signals are reliable, and that the main unresolved bottleneck is importance-map calibration and selectivity rather than quantizer form alone. Ongoing Stage3 experiments (multi-scale ROI heads on ResNet and DarkNet) are intended to close this gap through detector-facing operating-point analysis.
